% !PRINTED 464
\weberpage{464}
%<<CONT
mit $\overline{P}_2$ verbundenen Punkt der Curve $f$ gehen; denn sie
kann nicht durch die Punkte $P_7, P_3$ gehen, weil in diesen entweder
$\Delta$ oder $C$ verschwindet, sie kann aber auch nicht durch einen
Punkt der Grundcurve gehen, der kein Pol ist, weil sie sonst durch
die 168 verbundenen Punkte gehen müsste, und nicht durch den Pol
$\overline{P}_2$ gehen könnte. Dann kann man aber die Constante $h$
so bestimmen, dass $K^2 - hQ$ durch $f$ theilbar wird.

Der Quotient ist eine Invariante von niedrigerem als dem
$42^{\text{sten}}$, jedenfalls aber von geradem Grade, und wenn wir
ihn also durch die fundamentalen Invarianten darstellen, so kann
darin $K$ nicht vorkommen. Daraus folgt:

\begin{center}
\gesperrt{Die Invariante $K^2$ kann rational durch $f, \Delta, C$
ausgedrückt werden.}
\end{center}

Stellen wir nach diesem Satze $K^2$ in der Form dar:
\weq{3}{K^2 = \Sigma\, af^{\nu}\Delta^{\nu_1}C^{\nu_2},}
so können in dieser Summe, in der die $a$ numerische
Co\"efficienten sind, nur solche Glieder vorkommen, in denen
\[
2\nu + 3\nu_1 + 7\nu_2 = 21,
\]
und indem wir nun abzählen, welche Werthe von $\nu, \nu_1, \nu_2$
vorkommen können, erhalten wir das Resultat, dass zwischen den
Formen
\[
\begin{gathered}
K^2,\quad \Delta^7,\quad C^3,\quad f\Delta^4 C,\quad f^2\Delta C^2,\quad f^3\Delta^5,\\
f^4\Delta^2 C,\quad f^6\Delta^3,\quad f^7 C,\quad f^9\Delta
\end{gathered}
\]
eine lineare Relation mit numerischen Co\"efficienten besteht.
Stellen wir diese Relation in der Form dar:
\weq{4}{K^2 = \Phi\Delta + \Psi C,}
worin $\Phi, \Psi$ gleichfalls Invarianten sind, so können wir daraus
noch einen geometrischen Schluss ziehen.

Die Curven $\Delta$ und $C$ schneiden sich sicher nicht auf der Curve
$f$, weil die sämmtlichen Schnittpunkte von $\Delta$ mit $f$ die
$P_7$, die von $C$ mit $f$ die $P_3$ sind. Wenn aber $\Delta$ und
$C = 0$ sind, so ist auch $K = 0$, und folglich liegen
\gesperrt{alle Schnittpunkte von $\Delta$ und $C$ auf den 21 Hauptaxen der
Gruppe}\footnote{Die Relation (3) ist von \gesperrt{Gordan} durch die
Methoden der Invariantentheorie vollständig berechnet (Mathem.\
Annalen, Bd.\ XVII, S.\,371).

Wir wollen die dort gegebene Formel in
den von uns gebrauchten Zeichen hier angeben:
\[
K^2 = C^3 - 88f^2\Delta C^2 + 16.63f\Delta^4 C + 17.64f^4\Delta^2 C
- 256f^7 C + 27.64\Delta^7 - 128.469f^3\Delta^5 + 43.512f^6\Delta^3
- 2048f^9\Delta.
\]}.
